\chapter{Exemplos \LaTeX}
% exemplo de como inserir uma referencia adicional no sumario
\addcontentsline{toc}{chapter}{Exemplos que devem ser lidos :-)}

Esse capítulo tem exemplos de escrita utilizando o \LaTeX  utilizando \abnTeX, é muito simples escrever em \textbf{negrito}, \emph{italico}, ....

Leia o \ac{faq} do \abnTeX : \href{url}{https://github.com/abntex/abntex2/wiki/FAQ}

Existem diversos tutoriais para uso de \LaTeX, se você está utilizando esse modelo não precisará se preocupar com muitos dos detalhes técnicos do \LaTeX  e cuidar somente do seu texto.

Escolha seu editor : \href{url}{https://en.wikipedia.org/wiki/Comparison\_of\_TeX\_editors}


\section{Detalhes textuais}

O documento é dividido em capítulos, e cada capítulo dividido em seções utilizando o \abnTeX você pode dividir seus documentos nos níveis a seguir:

\begin{itemize}
\item chapter (1);
\item section (1.1);
\item subsection (1.1.1);
\item subsubsection (1.1.1.1);
\item subsubsubsection (1.1.1.1.1).
\end{itemize}

Tenha em mente que normalmente se utiliza no máximo o nível \emph{subsection}.
Ao definir as divisões do seu trabalho utilizando as diretivas do \LaTeX, elas são automaticamente inseridas no sumário do documento.

\subsection{Caracteres Reservados}



Alguns caracteres são reservados no \LaTeX e por isso para utilizar esses caracteres é necessario utilizar uma forma diferenciada de escrita. É possivel utilizar a macro \emph{symbol} com o codigo ascii do caracter desejado.
\begin{itemize}
\item barra invertida : \textbackslash   \symbol{92}    $\backslash$;
\item til  :  \symbol{126} ;
\item cifrão : \$;
\item sublinhado, \emph{underscore}, \emph{underline} : \_;
\item chaves : \} \{.
\end{itemize}



\subsection{Listas}

Em uma lista de itens cada item deve ser terminado por ponto e virgula, exceto o ultimo item que deve ter um ponto final.

\begin{itemize}
\item item 1;
\item item 2;
\item item ..;
\item item final.
\end{itemize}

\subsection{Elementos não textuais}

Elementos não textuais são aqueles que auxiliam o entendimento, não podem ficar "jogados" no texto, devem ser citados, cada elemento deve ser identificado por um \emph{label} único que permite a sua referencia, no texto utilizando \emph{ref} ou \emph{autoref}, esses elementos quando definidos corretamente também são inseridos nas listas presentes antes do sumário.


\subsection{Tabelas e Quadros}
A seção 3.32 da NBR14724:2011 define a Tabela como sendo uma "forma não discursiva de apresentar informações das quais o dado numérico se destaca como informação central" 

Quadros e tabelas são informações tabulares, mas Tabelas tem como objetivo apresentar números.

Uso de tabelas no \LaTeX : \href{url}{https://en.wikibooks.org/wiki/LaTeX/Tables}


\index{quadros}O \autoref{quadro-exemplo} é um exemplo de dados tabulares gerados em 
\LaTeX.

\begin{quadro}[htb]
\centering
\ABNTEXfontereduzida
\caption[Níveis de investigação]{Níveis de investigação.}
\label{quadro-exemplo}
\begin{tabular}{|p{2.6cm}|p{6.0cm}|p{2.25cm}|p{3.40cm}|}
  \hline
   \textbf{Nível de Investigação} & \textbf{Insumos}  & \textbf{Sistemas de Investigação}  & \textbf{Produtos}  \\
    \hline
    Meta-nível & Filosofia\index{filosofia} da Ciência  & Epistemologia &
    Paradigma  \\
    \hline
    Nível do objeto & Paradigmas do metanível e evidências do nível inferior &
    Ciência  & Teorias e modelos \\
    \hline
    Nível inferior & Modelos e métodos do nível do objeto e problemas do nível inferior & Prática & Solução de problemas  \\
   \hline
\end{tabular}
\legend{Fonte: Próprio Autor}
\end{quadro}



\index{tabelas}Já a \autoref{tab-exemplo} foi criada conforme o padrão do IBGE
requerido pelas normas da ABNT para documentos técnicos e acadêmicos. Números devem ser alinhados a direita.

\begin{table}[htb]
\centering
\caption{Um Exemplo de tabela}
\label{tab-exemplo}
\begin{tabular}{p{2.6cm}|r|r|r}
    \hline
   \textbf{Item} & \textbf{Janeiro}  & \textbf{Fevereiro}  & \textbf{Março}  \\
    \hline
    Classes & 2  & 10 & 20  \\
    \hline
    Linhas & 100  & 250 & 543 \\
    \hline
\end{tabular}
\fonte{Dados do Projeto}
\end{table}

\def\equationautorefname~#1\null{%
  Equação~(#1)\null
}


\index{equação}\index{Pitagoras}A \autoref{eq-pythagoras} mostra que também é possivel escrever equações diretamente em \LaTeX

\begin{equation}\label{eq-pythagoras}
a^2+b^2=c^2\,.
\end{equation}






% ---
\subsection{Figuras}
% ---

\index{figuras}Figuras podem ser criadas diretamente em \LaTeX,
como o exemplo da \autoref{fig_circulo}, ou inseridas a partir de arquivos externos como a \autoref{fig_logo}, que é o Logotipo do IFSP\index{IFSP}. \index{logotipo}

\begin{figure}[htb]
	\begin{center}
	    \setlength{\unitlength}{5cm}
		\begin{picture}(1,1)
		\put(0,0){\line(0,1){1}}
		\put(0,0){\line(1,0){1}}
		\put(0,0){\line(1,1){1}}
		\put(0,0){\line(1,2){.5}}
		\put(0,0){\line(1,3){.3333}}
		\put(0,0){\line(1,4){.25}}
		\put(0,0){\line(1,5){.2}}
		\put(0,0){\line(1,6){.1667}}
		\put(0,0){\line(2,1){1}}
		\put(0,0){\line(2,3){.6667}}
		\put(0,0){\line(2,5){.4}}
		\put(0,0){\line(3,1){1}}
		\put(0,0){\line(3,2){1}}
		\put(0,0){\line(3,4){.75}}
		\put(0,0){\line(3,5){.6}}
		\put(0,0){\line(4,1){1}}
		\put(0,0){\line(4,3){1}}
		\put(0,0){\line(4,5){.8}}
		\put(0,0){\line(5,1){1}}
		\put(0,0){\line(5,2){1}}
		\put(0,0){\line(5,3){1}}
		\put(0,0){\line(5,4){1}}
		\put(0,0){\line(5,6){.8333}}
		\put(0,0){\line(6,1){1}}
		\put(0,0){\line(6,5){1}}
		\end{picture}
	\end{center}
	\caption{\label{fig_circulo}A delimitação do espaço}
	\legend{Fonte: os autores}
\end{figure}


\begin{figure}[htb]
    \centering
	\includegraphics{\ifspprefixo/LogoUFSCar.jpg}
	\caption{\label{fig_logo}Logo IFSP}
	\legend{Fonte: IFSP}
\end{figure}




